\chapter{Normalization and Assumption Discharge}
\label{sec:p07-natural-deduction-normalization}

Two natural-deduction presentations can derive the same consequences
and behave differently under proof reduction. The nonvacuous presentation
on Sheet~\ref{sheet:nj} requires each discharged label to mark at least
one actual assumption occurrence. The standard presentation on
Sheet~\ref{sheet:p07-nj-vacuous} permits an empty discharged class.
This distinction is essential when a reduction deletes a branch.

\section{The discharge distinction}

\begin{hbdefinition}{def:p07-nj-presentations}{The two finite-tree calculi}
Write \(\mathsf{NJ}_{\mathrm{nv}}\) for the exact calculus of
Sheet~\ref{sheet:nj}, and \(\mathsf{NJ}_{\mathrm v}\) for
Sheet~\ref{sheet:p07-nj-vacuous}. They have the same introduction
and elimination rules for
\(\bot,\land,\lor,\to,\forall,\exists\), with defined negation,
total capture-avoiding term substitution, and the same eigenvariable
conditions. They differ only in allowing empty discharge classes for
\(\to I\), each minor derivation of \(\lor E\), and the minor
derivation of \(\exists E\) in \(\mathsf{NJ}_{\mathrm v}\).

Derivations are finite rooted trees. An open assumption is an actual
undischarged leaf occurrence, not an unused member of a supporting
context. Metalinguistic natural-number labels distinguish the classes discharged by particular rule
applications. Bound individual variables and discharge labels can be
renamed consistently; free individual variables remain fixed.
Formulas and conclusions are compared up to bound-variable
alpha-equivalence. Write \(\mathrm{OA}(\mathcal D)\) for the set
of formulas at the open leaves of a derivation \(\mathcal D\).
\end{hbdefinition}

\begin{hbtheorem}{thm:p07-nj-nonvacuous-obstruction}{proved=annex:p7.nj.nonvacuous-obstruction}{The nonvacuous restriction obstructs ordinary normalization}
There is a closed \(\mathsf{NJ}_{\mathrm{nv}}\) derivation of
\(A\to(B\to A)\), for distinct atomic \(A,B\), whose only
principal introduction--elimination detour cannot be contracted
within that calculus. It has no disjunction or existential
permutation available. Thus ordinary compatible principal and
\(\lor E/\exists E\) permutation conversions do not give a
normalization theorem for \(\mathsf{NJ}_{\mathrm{nv}}\).
The same derivation normalizes in \(\mathsf{NJ}_{\mathrm v}\).
\end{hbtheorem}

The issue is a rule condition, not the truth of the implication.
The following bridge preserves consequence while distinguishing
the permitted reduction behavior. The obstruction also appears in
Read's discussion of nonvacuous discharge
\cite[p.~18]{Read2014ProofTheoreticValidityManuscript}.

\begin{hbtheorem}{thm:p07-nj-vacuous-consequence-bridge}{proved=annex:p7.nj.consequence-bridge}{The same consequences, with different reduction behavior}
For every finite derivation in either presentation there is a
derivation in the other with the same conclusion and the same
open assumption occurrences, up to consistent bound-name and
label renaming. In particular, for every finite supporting
assumption set \(\Gamma\) and formula \(C\),
\[
 \Gamma\derives_{\mathsf{NJ}_{\mathrm{nv}}}C
 \quad\Longleftrightarrow\quad
 \Gamma\derives_{\mathsf{NJ}_{\mathrm v}}C.
\]
The translation into the nonvacuous system inserts a conjunction
introduction followed by projection whenever an empty discharge
class must be populated. It does not preserve normality.
\end{hbtheorem}

\section{Proof terms and the exact substitutions}

\begin{hbdefinition}{def:p07-nj-proof-terms}{Annotations for the full first-order tree language}
Use proof variables \(u,v,\ldots\) for labeled assumption classes,
and separate individual variables \(x,y,\ldots\) in first-order
terms and formula annotations. A proof term \(M:A\) records
one derivation of formula \(A\). The constructors are as follows:
\begin{enumerate}
\item \(u:A\) is an assumption; \(\langle M,N\rangle:A\land B\)
is conjunction introduction; \(\pi_1M:A\) and \(\pi_2M:B\)
are its two eliminations when \(M:A\land B\).
\item \(\iota_1^B M:A\lor B\) and \(\iota_2^A N:A\lor B\)
are the injections. The term
\(\mathsf{case}(M;u.N;v.P):C\) has \(M:A\lor B\),
branches \(N:C\) under \(u:A\) and \(P:C\) under \(v:B\).
The two proof variables are bound in their respective branches.
\item \(\lambda u^A.M:A\to B\), where \(M:B\), binds
the assumption class \(u\). Application \(MN:B\) has
\(M:A\to B\) and \(N:A\).
\item \(\mathsf{abort}_A(M):A\) has \(M:\bot\).
There is no introduction constructor for \(\bot\).
\item \(\Lambda x.M:\forall x A\) has \(M:A\) and
the registered universal eigenvariable condition. Its elimination
\(M[t]:A[t/x]\) has \(M:\forall x A\).
\item \(\mathsf{pack}_{x.A}(t,M):\exists x A\) has
\(M:A[t/x]\). The term
\(\mathsf{unpack}(M;y,u.N):C\) has \(M:\exists x A\)
and a branch \(N:C\) under \(u:A[y/x]\), with the registered
existential eigenvariable condition. The names \(y,u\) are
bound in this branch.
\end{enumerate}
The binding proof variables may have zero occurrences in
\(\mathsf{NJ}_{\mathrm v}\). Formula annotations and witness
terms are part of the syntax, even when an auxiliary termination
translation later ignores their individual-term detail.
\end{hbdefinition}

\begin{hbdefinition}{def:p07-nj-substitutions}{Proof substitution and individual substitution}
\(M[N/u]\) replaces all free occurrences of the proof variable
\(u\) by separately renamed copies of \(N\). Bound proof labels
are first made fresh for the free labels of the inserted proof,
and internal individual eigenvariables are made fresh for its
individual variables and open assumptions. The operation is
simultaneous: it does not substitute again inside an inserted copy.

\(M[t/x]\) applies capture-avoiding individual-term substitution
to every witness term and formula annotation, with consistent
changes to the derivation's individual-variable occurrences.
It does not replace proof assumptions by proofs. Internal
\(\forall I\) and \(\exists E\) eigenvariables are first
renamed away from \(x\) and the variables of \(t\).
\end{hbdefinition}

\begin{hblemma}{thm:p07-nj-substitution}{proved=annex:p7.nj.substitution}{Substitution preserves the standard derivation conditions}
Suppose \(M:C\) is a \(\mathsf{NJ}_{\mathrm v}\) derivation
with an open class \(u:A\), and \(N:A\) is another such
derivation. Then \(M[N/u]\) is a derivation of \(C\), with
open assumptions contained in those of \(M\) other than class
\(u\), together with the open assumptions of \(N\).
No new assumption formula is introduced. Individual substitution
produces a derivation of the correspondingly substituted formula
from the correspondingly substituted open assumptions.
Both assertions include all first-order quantifier cases and
the possibility of zero, one or many inserted copies.
\end{hblemma}

\section{Principal conversions}

\begin{hbdefinition}{def:p07-nj-principal-conversions}{The complete introduction--elimination table}
With the constructor types and freshness conditions above, the
principal conversions are
\[
\begin{aligned}
 \pi_1\langle M,N\rangle&\longrightarrow M,&
 \pi_2\langle M,N\rangle&\longrightarrow N,\\
 (\lambda u^A.M)N&\longrightarrow M[N/u],\\
 \mathsf{case}(\iota_1^B M;u.N;v.P)&\longrightarrow N[M/u],\\
 \mathsf{case}(\iota_2^A M;u.N;v.P)&\longrightarrow P[M/v],\\
 (\Lambda x.M)[t]&\longrightarrow M[t/x],\\
 \mathsf{unpack}(\mathsf{pack}_{x.A}(t,M);y,u.N)
       &\longrightarrow N[t/y][M/u].
\end{aligned}
\]
In the last line, \(y\) is first fresh for \(t,M\), and
the substituted class has formula \(A[t/x]\). There is no
principal \(\bot I/\bot E\) pair, because the calculus has
no \(\bot I\). The conversions operate on one selected
occurrence inside a proof tree; compatible closure permits
that occurrence at any proof-term position.
\end{hbdefinition}

\begin{hbtheorem}{thm:p07-nj-principal-preservation}{proved=annex:p7.nj.principal-preservation}{Every principal conversion preserves conclusion and allowed assumptions}
Each displayed principal conversion takes a well-formed
\(\mathsf{NJ}_{\mathrm v}\) derivation to one with the same
conclusion and no new open assumption formula, including when
the conversion occurs inside a larger derivation. The descendant
discharge classes that become empty remain legal in this presentation.
\end{hbtheorem}

\section{Permutation and absurdity conversions}

\begin{hbdefinition}{def:p07-nj-elimination-frames}{Every permitted following elimination}
An elimination frame \(E[-]\) has its one hole in the major-premise
position and consists of exactly one of the following seven forms.
Its auxiliary proofs, witness term and final formula are fixed when
the hole is filled:
\[
\begin{array}{c|l}
\text{rule}&E[-]\\ \hline
\land E_1&\pi_1[-]\\
\land E_2&\pi_2[-]\\
\to E&[-]L\\
\lor E&\mathsf{case}([-];c.L;d.R)\\
\forall E&[-][t]\\
\exists E&\mathsf{unpack}([-];z,c.L)\\
\bot E&\mathsf{abort}_B([-])
\end{array}
\]
For each of these seven rows there are two permutation conversions:
\[
\begin{aligned}
 E[\mathsf{case}(M;u.N;v.P)]
 &\longrightarrow
   \mathsf{case}(M;u.E_1[N];v.E_2[P]),\\
 E[\mathsf{unpack}(M;y,u.N)]
 &\longrightarrow
   \mathsf{unpack}(M;y,u.E[N]).
\end{aligned}
\]
Here \(E_1,E_2\) are separately freshened copies of the entire
frame and its auxiliary derivations, not just copies of its last
rule symbol. This specifies all fourteen permutations, including
a case after a case, an unpacking after a case, a case after an
unpacking, and nested unpackings.

Before enlarging a branch scope, rename its \(u,v\), or \(y,u\),
fresh for the following frame, its auxiliary proofs, every open
assumption there, its individual terms and its final formula.
For the universal-elimination row this includes the term \(t\).
For the existential-elimination row the following eigenvariable
\(z\) must also be fresh for the earlier branch assumptions,
which were already discharged on the left and need not previously
have belonged to its open context. The two case branches receive
distinct copied \(z\)'s and bound proof labels. Nested unpackings
use distinct eigenvariables and distinguished proof classes.
Free external individual variables are not renamed by these operations.
\end{hbdefinition}

\begin{hbdefinition}{def:p07-nj-bottom-conversions}{Absorption of an elimination after absurdity}
For each of the seven frames above, if its major premise has
formula \(A\) and its final conclusion formula \(B\), use
\[
                  E[\mathsf{abort}_A(M)]
                   \longrightarrow \mathsf{abort}_B(M).
\]
Thus the two projections return an absurdity elimination to the
selected conjunct; application returns one to the consequent;
case and unpacking return one to their branch conclusion;
universal elimination returns one to \(A_0[t/x]\) when its
major formula is \(\forall x A_0\); and an absurdity elimination
after \(\mathsf{abort}_\bot(M)\) returns the single final
absurdity elimination. All auxiliary proofs of the following
frame disappear. The retained \(M\) still concludes \(\bot\).
\end{hbdefinition}

\begin{hbtheorem}{thm:p07-nj-permutation-preservation}{proved=annex:p7.nj.permutation-preservation}{All fourteen permutations and seven absorptions preserve the derivation}
With the stated freshness and copying conditions, every displayed
permutation and absurdity absorption preserves the conclusion
and introduces no new open assumption formula. The conclusion
and external free variables are retained, including the cases
with quantified formulas, term instantiations and nested
existential eliminations. The same holds in every proof context.
\end{hbtheorem}

\begin{hbdefinition}{def:p07-nj-vacuous-normality}{The complete normality condition and reduction relation}
Let \(\to_{\mathrm D}\), \(\to_{\mathrm P}\), and
\(\to_{\mathrm B}\) be the compatible closures of the seven
principal, fourteen permutation, and seven absurdity-absorption
rules respectively. Alpha-renaming is equality of representations,
not an extra directed step. A proof is \emph{normal} when no
occurrence anywhere admits any of these three kinds of step.
Equivalently, an elimination's major premise is not the result
of its matching introduction, of a case/unpacking, or of an
absurdity elimination. No eta expansion, permutation from a
minor-premise position, or commuting introduction is included.
\end{hbdefinition}

\section{An explicit termination argument}

The proof uses the continuation translation and arithmetic measure
of de Groote's full first-order treatment, with all actual
quantifier annotations retained on the derivation
\cite[Definitions~3, 6--10, 13--16, 19--20]{DeGroote2002NJPermutation}.
The auxiliary arrow-only normalization theorem used below is
\cref{thm:p07-stlc-strong-normalization}, whose full proof includes
open terms and finite beta-reduction heights.

\begin{hbdefinition}{def:p07-nj-erasure}{Erasing individual detail only in the termination simulation}
Write \(e(M)\) for the proof tree with formula annotations and
individual terms erased. Propositional constructors and proof-variable
binding remain. For the quantified constructors set
\[
\begin{aligned}
e(\Lambda x.M)&=\mathsf{ai}(e(M)),&
e(M[t])&=\mathsf{ae}(e(M)),\\
e(\mathsf{pack}_{x.A}(t,M))&=\mathsf{pack}(e(M)),&
e(\mathsf{unpack}(M;y,u.N))&=
        \mathsf{unpack}(e(M);u.e(N)).
\end{aligned}
\]
Also erase the formula subscript of \(\mathsf{abort}\) and the
missing-summand annotation of an injection. This is a simulation
map; it does not identify actual derivations with different witnesses.
In particular, an actual universal principal step still erases
to \(\mathsf{ae}(\mathsf{ai}(U))\to U\), removing two proof
constructors. Write \(U,V,W\) for erased terms below.
\end{hbdefinition}

\begin{hbdefinition}{def:p07-nj-cps-types}{Arrow-only types for the simulation}
Fix distinct simple base types \(o,\iota\), and a base type
\(p_P\) for each predicate symbol \(P\). Define value and
computation types by \(A^{\mathrm c}=(A^{\mathrm v}\to o)\to o\)
and the following recursion, with arrows associating to the right:
\[
\begin{array}{c|l}
A&A^{\mathrm v}\\ \hline
P(t_1,\ldots,t_n)&p_P\\
\bot&o\\
B\to C&B^{\mathrm c}\to C^{\mathrm c}\\
B\land C&(B^{\mathrm c}\to C^{\mathrm c}\to o)\to o\\
B\lor C&(B^{\mathrm c}\to o)\to(C^{\mathrm c}\to o)\to o\\
\forall x B&\iota\to B^{\mathrm c}\\
\exists x B&(\iota\to B^{\mathrm c}\to o)\to o
\end{array}
\]
These arrows are simple-type constructors in the target calculus,
not additional object-language implication rules. Individual
substitution does not change either translated type. Translate
each open proof variable \(u:A\) as \(u:A^{\mathrm c}\), and
add a fresh target variable \(\star:\iota\). This variable
is a marker in an open STLC context; it is not an individual
witness postulated in the original logic.
\end{hbdefinition}

\begin{hbdefinition}{def:p07-nj-modified-cps}{The complete modified continuation translation}
Define \(\mathcal T(U)=\lambda k.\mathcal C(U,k)\), with
\(k\) fresh. For a target continuation term \(K\), define
\(\mathcal C\) by the following simultaneous structural recursion.
Every displayed target binder is fresh for its operands; all
target lambdas carry the simple type determined by the preceding
type table and the typing theorem below.
\[
\begin{array}{c|l}
U&\mathcal C(U,K)\\ \hline
u&uK\\
\lambda u.V&K(\lambda u.\mathcal T(V))\\
VW&\mathcal C(V,\lambda f.f\mathcal T(W)K)\\
\langle V,W\rangle&K(\lambda p.p\mathcal T(V)\mathcal T(W))\\
\pi_1V&\mathcal C(V,\lambda p.p(\lambda u.\lambda v.uK))\\
\pi_2V&\mathcal C(V,\lambda p.p(\lambda u.\lambda v.vK))\\
\iota_1V&K(\lambda u.\lambda v.u\mathcal T(V))\\
\iota_2V&K(\lambda u.\lambda v.v\mathcal T(V))\\
\mathsf{case}(V;u.W;v.Z)&
 \mathcal C(V,\lambda s.s(\lambda u.\mathcal C(W,K))
                               (\lambda v.\mathcal C(Z,K)))\\
\mathsf{ai}(V)&K(\lambda\xi.\mathcal T(V))\\
\mathsf{ae}(V)&\mathcal C(V,\lambda f.f\star K)\\
\mathsf{pack}(V)&K(\lambda p.p\star\mathcal T(V))\\
\mathsf{unpack}(V;u.W)&
 \mathcal C(V,\lambda s.s(\lambda\xi.\lambda u.\mathcal C(W,K)))\\
\mathsf{abort}(V)&\mathcal C(V,(\lambda z.\lambda r.r)K)
\end{array}
\]
The dummy \(\xi\)'s have type \(\iota\). In the last row
\(z\) has the type of \(K\), and \(r:o\). Keep that
administrative beta redex uncontracted in the definition: its
argument preserves occurrences of \(K\) needed for strict
simulation of reductions inside auxiliary derivations.
\end{hbdefinition}

\begin{hbtheorem}{thm:p07-nj-cps-simulation}{proved=annex:p7.nj.cps-simulation}{Typing and the two simulation properties}
If \(M:A\) is a \(\mathsf{NJ}_{\mathrm v}\) derivation, then
\(\mathcal T(e(M)):A^{\mathrm c}\) is an open arrow-only STLC
term in its translated context with \(\star:\iota\).
Every actual principal step \(M\to_{\mathrm D}N\) gives
\[
             \mathcal T(e(M))\to_\beta^+
                          \mathcal T(e(N)),
\]
where \(+\) means at least one beta step. Every actual
permutation \(M\to_{\mathrm P}N\) gives alpha-equivalent
target terms. Both statements hold under all proof contexts,
not only at the root.
\end{hbtheorem}

\begin{hbdefinition}{def:p07-nj-permutation-measure}{A measure that accounts for copied branches}
Define positive integers \(n(U),w(U)\) on erased proof terms:
\[
\begin{array}{c|c|l}
U&n(U)&w(U)\\ \hline
u&1&1\\
\lambda u.V&1&w(V)\\
VW&n(V)&w(V)+n(V)w(W)\\
\langle V,W\rangle&1&w(V)+w(W)\\
\pi_iV&n(V)&w(V)+n(V)\\
\iota_iV&1&w(V)\\
\mathsf{case}(V;u.W;v.Z)&2n(V)(n(W)+n(Z))&
                 w(V)+n(V)(w(W)+w(Z))\\
\mathsf{ai}(V)&1&w(V)\\
\mathsf{ae}(V)&n(V)&w(V)+n(V)\\
\mathsf{pack}(V)&1&w(V)\\
\mathsf{unpack}(V;u.W)&2n(V)n(W)&w(V)+n(V)w(W)\\
\mathsf{abort}(V)&n(V)&w(V)+n(V)
\end{array}
\qquad (i=1,2).
\]
For the erased version of each elimination frame there are
positive integers \(a_E,b_E\), independent of the hole, given by
\[
\begin{array}{c|c|c}
E[-]&a_E&b_E\\ \hline
[-]L&w(L)&1\\
\pi_1[-],\ \pi_2[-],\ \mathsf{ae}([-]),\ \mathsf{abort}([-])&1&1\\
\mathsf{case}([-];c.L;d.R)&w(L)+w(R)&2(n(L)+n(R))\\
\mathsf{unpack}([-];c.L)&w(L)&2n(L)
\end{array}
\]
They satisfy \(w(E[U])=w(U)+n(U)a_E\) and
\(n(E[U])=n(U)b_E\).
\end{hbdefinition}

\begin{hblemma}{thm:p07-nj-permutation-decrease}{proved=annex:p7.nj.permutation-decrease}{Strict decrease for every permutation}
Every compatible permutation preserves \(n\) and strictly
decreases \(w\). At a case permutation the decrease is exactly
\(n(M)(n(N)+n(P))a_E\); at an unpacking permutation it is
\(n(M)n(N)a_E\), in the erased notation of their defining
patterns. These quantities are positive even when the following
elimination copies entire branch derivations.
\end{hblemma}

\begin{hbtheorem}{thm:p07-nj-vacuous-normalization}{proved=annex:p7.nj.vacuous-normalization}{Full first-order normalization with empty discharge classes}
Every finite \(\mathsf{NJ}_{\mathrm v}\) derivation admits a
finite sequence of the declared principal, permutation and
absurdity conversions ending at a normal derivation with the
same conclusion and an open-assumption formula set contained
in the original one, up to bound-variable alpha-equivalence.
The theorem includes all six logical constructors and the
exact discharge and eigenvariable conditions on its own sheet.

A terminating strategy first contracts principal or permutation
redexes until none remain, then contracts absurdity-absorption
redexes until none remain. No claim about arbitrary interleavings
of all three classes is required by this theorem.
\end{hbtheorem}

\begin{hbcorollary}{thm:p07-nj-principal-permutation-normalization}{proved=annex:p7.nj.principal-permutation-normalization}{Normalization using only principal and permutation conversions}
Every finite \(\mathsf{NJ}_{\mathrm v}\) derivation has a
derivation with the same conclusion and no new open assumption
formula, reachable using only \(\to_{\mathrm D}\) and
\(\to_{\mathrm P}\), which contains neither of those redex
classes. This is the normality condition that excludes
introduction--elimination detours and the specified
\(\lor E/\exists E\) permutation detours but does not
add a condition on absurdity-elimination chains.
\end{hbcorollary}

\begin{hbmethod}{def:p07-nj-normalization-use}{Which conclusion the normalization supplies}
Use the formula-annotated proof tree and its original assumption
ownership throughout the actual reductions. The erasure and
continuation translation certify termination; they do not
replace the quantified derivation that is returned. A proof
from the nonvacuous presentation may first be regarded as a
proof in \(\mathsf{NJ}_{\mathrm v}\) and normalized there.
Translating it back with the consequence bridge can insert
new detours, so it does not give a nonvacuous normal form.

The classical maximal-formula/segment presentations in
\cite{Prawitz1965,TroelstraSchwichtenberg2000} and the
continuation proof used here are different proof routes.
Theorem~\ref{thm:p07-nj-vacuous-normalization} supplies the
full stated first-order result by the latter route. It does
not identify a source's arrow-only strong-normalization proof
with a proof for quantified disjunction and existential elimination.
\index{normalization!empty discharge classes}
\index{natural deduction!nonvacuous variant}
\index{permutation conversion!quantified natural deduction}
\end{hbmethod}
